%! Introducing Statistics: Why Be Normal? — Unit 6, Lesson 6.1 — Guided Notes
\documentclass[10pt]{article}
\usepackage{apstats-article}
\usepackage{apstats-boxes}
\newcommand{\TallMath}[1]{$\displaystyle #1\rule[-1.4em]{0pt}{3.2em}$}

\begin{document}

\pageheader{Unit 6, Lesson 6.1}{Guided Notes}
\namedateperiod

\vspace{0.15in}

\begin{objectivebox}
By the end of this lesson, I will be able to:
\begin{enumerate}
  \item \writeline
  \item \writeline
\end{enumerate}
\end{objectivebox}

\vspace{0.1in}

\begin{vocabbox}
\termblanklong{Statistical inference}
\termblanklong{Confidence interval}
\termblanklong{Significance test (hypothesis test)}
\termblanklong{Null hypothesis ($H_0$)}
\termblanklong{Alternative hypothesis ($H_a$)}
\end{vocabbox}

\vspace{0.1in}

\begin{notesbox}{The Big Shift: From Unit 5 to Unit 6}
\textbf{Unit 5 question:} Given $p = 0.60$ and $n = 100$, what is the probability that
$\hat p > 0.65$?

\textbf{Unit 6 question:} Given $\hat p = 0.65$ and $n = 100$, what can we say about
the \emph{unknown} population proportion $p$?

The key: since the sampling distribution of $\hat p$ is approximately
\blank{3.5cm} when conditions are met, we can \blank{10cm}.

\end{notesbox}

\vspace{0.1in}

\begin{notesbox}{Two Types of Inference}
\renewcommand{\arraystretch}{1.8}
\begin{tabular}{@{} >{\bfseries}p{0.20\linewidth} p{0.36\linewidth} p{0.36\linewidth} @{}}
  & \textbf{Confidence Interval} & \textbf{Significance Test} \\
  \hline
  Goal    & \writeline & \writeline \\
  Question asked & ``Where is \blank{1.5cm}?'' & ``Is \blank{2cm} consistent with my data?'' \\
  Answer gives & A range of \blank{2.5cm} values & A decision: reject or \blank{3cm} $H_0$ \\
\end{tabular}
\end{notesbox}

\vspace{0.1in}

\begin{notesbox}{Conditions for Using a Normal Model}
Before using any inference procedure in Unit 6, verify:

\begin{enumerate}
  \item \textbf{Random:} The data come from \blank{8cm}.
  \item \textbf{Large Counts:} Both $n\hat{p} \ge$ \blank{1cm} and $n(1-\hat{p}) \ge$ \blank{1cm}.
  \item \textbf{Independence (10\% rule):} If sampling without replacement, $n \le$ \blank{4cm}.
\end{enumerate}

When all three conditions are met, the sampling distribution of $\hat p$ is approximately
\blank{8cm} with mean \blank{2.5cm} and standard deviation \blank{4cm}.
\end{notesbox}

\vspace{0.1in}

\begin{practicebox}
\textbf{Quick Check.} Classify each research question as calling for a \emph{confidence interval}
(CI) or a \emph{significance test} (Test).

\begin{enumerate}
  \item ``What proportion of U.S. adults exercise daily?'' \hfill \blank{2.5cm}
  \item ``Is there evidence that more than half of students prefer online homework?'' \hfill \blank{2.5cm}
  \item ``We want to estimate, within 3 percentage points, the rate of vaccine hesitancy.''
        \hfill \blank{2.5cm}
\end{enumerate}
\end{practicebox}

\end{document}
